\documentclass{article}
\usepackage[UTF8]{ctex}
\usepackage{amsmath,amssymb,amsfonts}
\usepackage{enumitem}
\usepackage{geometry}
\geometry{a4paper, margin=1in}
% === 错题本排版优化（不参与题目标识） ===
\setlength{\parindent}{0pt}
\setlength{\parskip}{0.35em plus 0.10em minus 0.08em}
\setlength{\emergencystretch}{3em}
\setlist{itemsep=0.25em, topsep=0.40em, parsep=0pt}
\allowdisplaybreaks[2]
% === 错题本排版优化结束 ===
\begin{document}
\section*{长安大学2026年高等代数真题}
\section*{176．长安大学2026年高等代数试卷}
一．填空题．每题5分，共30分．\\

\subsection{题目1}
1．设 \(f(x)\) 是一个实系数三次多项式，满足 \((x+1)^{2}\) 整除 \(f(x)+1,(x-1)^{2}\) 整除 \(f(x)-1\) ，则 \(f(x)=\)

\subsection{题目2}
2．实数域上四元二次型的不同规范形的个数是 \(\_\_\_\_\)．

\subsection{题目3}
3．设 \(A\) 为 3 阶方阵，且 \(|A|=\frac{1}{2}\) ，则 \(\left|(3 A)^{-1}-2 A^{*}\right|=\) \(\_\_\_\_\)．

\subsection{题目4}
4．设 \(A\) 为行列式等于1的3阶正交矩阵，则 \(A\) 必有实特征值 \(\_\_\_\_\)．

\subsection{题目5}
5．所有与 \(n\) 阶方阵 \(A\) 可交换的矩阵集合 \(C(A)\) 是 \(P^{n \times n}\) 的一个子空间，当对角矩阵

\[
A=\operatorname{diag}\left\{a_{1}, a_{2}, \cdots, a_{n}\right\}
\]

中对角元素 \(a_{i}(1 \leq i \leq n)\) 互不相同时，\(C(A)\) 的维数为 \(\_\_\_\_\)．

\subsection{题目6}
6．\(n\) 阶方阵 \(A\) 满足 \(A^{3}=2 E, B=A^{2}-2 A+2 E\) ，则 \(B^{-1}=\) \(\_\_\_\_\)．

二．（10 分）设 \(f(x)\) 是有理数域上的多项式，证明：\(f(x)\) 在有理数域上不可约当且仅当对任意有理数\(a, b\) ，其中 \(a \neq 0\) ，多项式 \(g(x)=f(a x+b)\) 在有理数域上不可约．

三．（10 分）设 \(n \geq 2\) ，计算 \(n\) 阶行列式

\[
D_{n}=\left|\begin{array}{cccc}
\sin \left(2 \alpha_{1}\right) & \sin \left(\alpha_{1}+\alpha_{2}\right) & \cdots & \sin \left(\alpha_{1}+\alpha_{n}\right) \\
\sin \left(\alpha_{2}+\alpha_{1}\right) & \sin \left(2 \alpha_{2}\right) & \cdots & \sin \left(\alpha_{2}+\alpha_{n}\right) \\
\vdots & \vdots & & \vdots \\
\sin \left(\alpha_{n}+\alpha_{1}\right) & \sin \left(\alpha_{n}+\alpha_{2}\right) & \cdots & \sin \left(2 \alpha_{n}\right)
\end{array}\right|
\]

四．（15 分）设非齐次线性方程组

\[
\left\{\begin{array}{c}
a_{11} x_{1}+a_{12} x_{2}+\cdots+a_{1 \mathrm{n}} x_{\mathrm{n}}=b_{1}, \\
a_{21} x_{1}+a_{22} x_{2}+\cdots+a_{2 \mathrm{n}} x_{\mathrm{n}}=b_{2}, \\
\cdots \cdots \\
a_{\mathrm{s} 1} x_{1}+a_{\mathrm{s} 2} x_{2}+\cdots+a_{\mathrm{sn}} x_{\mathrm{n}}=b_{\mathrm{s}}
\end{array}\right.
\]

的系数矩阵为 \(A\) ，向量 \(\beta=\left(b_{1}, b_{2}, \cdots, b_{s}\right)^{\prime}\) 。\\
（1）若方程组有解且秩 \((A)=r\) ，求方程组的解向量中最多线性无关解的数目。\\
（2）若方程组对任意的 \(s\) 维列向量 \(\beta\) 都有解，求秩 \((A)\) 。

五．（15 分）设

\[
\alpha=\left(\begin{array}{l}
1 \\
2 \\
3 \\
4
\end{array}\right), A=\left(\begin{array}{cccc}
2 & 1 & 3 & 4 \\
2 & 5 & 5 & 8 \\
3 & 6 & 10 & 11 \\
4 & 8 & 12 & 17
\end{array}\right)-\frac{1}{30^{2025}}\left(\alpha \alpha^{\prime}\right)^{2026}, B=\left(\begin{array}{cccc}
2 & 1 & 3 & 4 \\
0 & 2 & 1 & 3 \\
0 & 0 & 2 & 1 \\
0 & 0 & 0 & 2
\end{array}\right)
\]

（1）求矩阵 \(A\) 。\\
（2）若 \(X\) 满足 \(X\left(E-B^{-1} A\right)^{\prime} B^{\prime}=E\) ，求矩阵 \(X\) 。

六．（15 分）已知矩阵 \(A=\left(\begin{array}{ccc}2 & -1 & -1 \\ 2 & -1 & -2 \\ -1 & 1 & 2\end{array}\right)\) ，求 \(A^{100}\) ．

七．（15 分）设 \(A\) 为 \(n\) 阶方阵，\(A\) 的每一行每一列恰好有一个元素等于1，其余元素均为0，证明：存在正整数 \(k\) ，使得 \(A^{k}=E\) 。

八．（15 分）设 \(\sigma\) 是数域 \(P\) 上线性空间 \(V\) 上的线性变换，且 \(\sigma^{2}=\sigma\) ．证明：

\begin{itemize}
  \item[（1）] \(\sigma\) 的特征值为1或0。
  \item[（2）] \(V=\sigma^{-1}(0) \oplus \sigma(V)\) ．
\end{itemize}

九．（15 分）求矩阵

\[
A=\left(\begin{array}{ccc}
-4 & 9 & -4 \\
-9 & 18 & -8 \\
-15 & 29 & -13
\end{array}\right)
\]

的不变因子，初等因子，最小多项式，若尔当标准形与有理标准形．

十．（10 分）设 \(A, B\) 为 \(m \times n\) 实矩阵，秩 \((A+B)=n\) 。证明：\(A^{\prime} A+B^{\prime} B\) 为正定矩阵。
\end{document}
